\documentclass{beamer}
\usetheme{Madrid}
\usecolortheme{default}

\usepackage{amsmath}
\usepackage{amssymb}
\usepackage{bm}
\usepackage{graphicx}
\usepackage{listings}
\usepackage{xcolor}

% Configure Python code formatting
\lstset{
    language=Python,
    basicstyle=\ttfamily\tiny,
    keywordstyle=\color{blue},
    stringstyle=\color{green!50!black},
    commentstyle=\color{red!80!black},
    showstringspaces=false,
    frame=single,
    breaklines=true
}

\title{LTI System Frequency and Step Response}
\subtitle{Problem 1.1.19 (EC 2007)}
\author{ee25btech11019}
\date{}

\begin{document}

\begin{frame}
    \titlepage
\end{frame}

\begin{frame}{Problem Statement}
    The frequency response of a linear time-invariant system is given by $H(f) = \frac{5}{1 + j10\pi f}$.
    
    The step response of the system is:
    
    \vspace{0.5cm}
    a) $5(1 - e^{-5t})u(t)$
    
    b) $5(1 - e^{-t/5})u(t)$
    
    c) $\frac{1}{5}(1 - e^{-5t})u(t)$
    
    d) $\frac{1}{5}(1 - e^{-t/5})u(t)$
\end{frame}

\begin{frame}{Transformation to the s-Domain}
    We first map the system from the Fourier frequency domain to the Laplace $s$-domain. 
    
    The relationship between the continuous Fourier transform and the Laplace transform (evaluating along the imaginary axis) is:
    \begin{equation}
        s = j\omega = j2\pi f
    \end{equation}
    
    We rewrite the denominator of $H(f)$ to isolate this term:
    \begin{equation}
        j10\pi f = 5(j2\pi f) = 5s
    \end{equation}
    
    Substituting this into the frequency response gives the transfer function $H(s)$:
    \begin{equation}
        H(s) = \frac{5}{1 + 5s} = \frac{1}{s + 0.2}
    \end{equation}
\end{frame}

\begin{frame}{Deriving the Step Response}
    The step response $Y(s)$ in the Laplace domain is the product of the system transfer function $H(s)$ and the step input $X(s)$.
    
    The Laplace transform of the Heaviside step function $u(t)$ is $X(s) = \frac{1}{s}$.
    
    \begin{equation}
        Y(s) = H(s)X(s) = \left( \frac{1}{s + 0.2} \right) \left( \frac{1}{s} \right) = \frac{1}{s(s + 0.2)}
    \end{equation}
    
    We perform a partial fraction expansion to prepare for the inverse transform:
    \begin{equation}
        \frac{1}{s(s + 0.2)} = \frac{A}{s} + \frac{B}{s + 0.2}
    \end{equation}
    
    Solving for the coefficients yields $A = 5$ and $B = -5$:
    \begin{equation}
        Y(s) = \frac{5}{s} - \frac{5}{s + 0.2}
    \end{equation}
\end{frame}

\begin{frame}{Inverse Laplace Transform}
    We now apply the standard inverse Laplace transforms to revert to the time domain:
    \begin{align}
        \mathcal{L}^{-1} \left\{ \frac{1}{s} \right\} &= u(t) \\
        \mathcal{L}^{-1} \left\{ \frac{1}{s + a} \right\} &= e^{-at}u(t)
    \end{align}
    
    Applying these to our expanded equation:
    \begin{equation}
        y(t) = 5u(t) - 5e^{-0.2t}u(t)
    \end{equation}
    \begin{equation}
        y(t) = 5(1 - e^{-t/5})u(t)
    \end{equation}
    
    \vspace{0.5cm}
    \textbf{Correct Option: (b)}
\end{frame}

\begin{frame}{System Physics: Deriving the Differential Equation}
    While the Laplace transform provides an algebraic pathway, we must rigorously define the governing Differential Equation of the system.
    
    Starting from our derived transfer function:
    \begin{equation}
        H(s) = \frac{Y(s)}{X(s)} = \frac{5}{5s + 1}
    \end{equation}
    
    We cross-multiply to relate the complex output $Y(s)$ to the input $X(s)$:
    \begin{equation}
        Y(s)(5s + 1) = 5X(s)
    \end{equation}
    \begin{equation}
        5sY(s) + Y(s) = 5X(s)
    \end{equation}
\end{frame}

\begin{frame}{Time-Domain Differential Equation}
    We apply the Inverse Laplace Transform. By the derivative property, multiplying by $s$ in the complex domain equates to a time derivative $\frac{d}{dt}$ in the time domain (assuming zero initial conditions):
    \begin{equation}
        5 \frac{dy(t)}{dt} + y(t) = 5x(t)
    \end{equation}
    
    To find the step response, our input $x(t)$ is the Heaviside step function $u(t)$:
    \begin{equation}
        5 \frac{dy(t)}{dt} + y(t) = 5u(t)
    \end{equation}
    
    This derives the exact first-order linear ordinary differential equation (ODE) of the system. In standard form $\tau y' + y = Kx$, we clearly see the system time constant is $\tau = 5$, which fundamentally proves why the decay envelope of our solution must be $e^{-t/5}$.
\end{frame}

\begin{frame}{Computational Verification Plot}
    \begin{figure}[H]
        \centering
        % Ensure you run the Python script first to generate this image
        \includegraphics[width=0.85\textwidth]{lti_response.png}
        \caption{Perfect overlap between the SciPy numerical simulation and SymPy's derived analytical equation.}
    \end{figure}
\end{frame}

\begin{frame}[fragile]{Python Verification}
\begin{lstlisting}
import sympy as sp
from scipy import signal
import matplotlib.pyplot as plt
import numpy as np

# 1. SYMBOLIC ALGEBRA ENGINE (SymPy)
s, t = sp.symbols('s t')
H_s = 5 / (1 + 5*s)
X_s = 1 / s # Step input in Laplace domain
Y_s = H_s * X_s
y_t = sp.inverse_laplace_transform(Y_s, s, t) # Exact analytical solve
print(f"SymPy Analytical Output: {y_t}")

# 2. NUMERICAL SIGNAL ENGINE (SciPy)
# Transfer function H(s) = 5 / (5s + 1) -> num=[5], den=[5, 1]
sys = signal.TransferFunction([5], [5, 1])
time, response = signal.step(sys)

# 3. VERIFICATION PLOT
plt.figure(figsize=(7, 4))
plt.plot(time, response, 'b-', lw=3, label='SciPy Numerical Step Response')
t_vals = np.linspace(0, max(time), 100)
y_sym = 5 * (1 - np.exp(-t_vals/5))
plt.plot(t_vals, y_sym, 'r--', lw=2, label=r'SymPy Analytical: $5(1 - e^{-t/5})$')

plt.title("LTI System: Symbolic vs Numerical Verification")
plt.xlabel("Time (t)")
plt.ylabel("Amplitude y(t)")
plt.legend(loc="lower right")
plt.grid(True)
plt.savefig('lti_response.png', dpi=150)
\end{lstlisting}
\end{frame}


\end{document}
